\scanpage{54}
Symétrie du cube : c'est l'homothétie $\unsure{\varphi|q}\unsure{\mapsto|-}-\varphi$. Du pt de \unsure{vue|vu} combinatoire, c'est l'autom $I'\to\sigma I'$ dans l'ens des ``sections partielles'' de $B$ sur $I$. Du pt de \unsure{vue|vu} maquette du cube, c'est l'autom $\sigma$ de $(B,\sigma)$. C'est l'unique autom \inserted{$u$} ayant la

\deleted{\uline{Prop} Soient $F,F'$ deux facettes \inserted{fermées} non vides, Alors $F\cap F'\ne\varnothing$ \uline{ou} $F\cap\sigma F'\ne\varnothing$. (Ceci caractérise l'autom $\sigma$ ?) \deleted{Les facettes} Celà revient à ceci : Si $I',I''\subset B$ sont deux sections partielles de $B$ sur $I=B/\sigma$, alors ou bien}

propriété : si $F$ est une facette de codim un (une hyperfacette), alors $\sigma F$ est une hyperfacette ne \deleted{\unsure{se}} rencontrant pas $F$. Plus gén, pour \deleted{\unsure{deux}} \inserted{tte} facette \inserted{$F$} \deleted{de dim \unsure{$d$}} à part celle de dim $n$, $F$ et $\sigma F$ ne se rencontrent pas (en fait, les hyperplans d'appui ne se \deleted{\unsure{rencontrent}} pas et ils sont $\parallel$)

Aussi : le centre de $\operatorname{Aut}(C)$ est le groupe $\{1,\sigma\}$ [\uline{NB} si $n\geq3$, le centre de $\unsure{\mathfrak S_I|\mathscr S_I}$ est $\{1\}$ \ldots], donc $\sigma$ est l'unique élément central $\ne1$.

% slanted
\margin{\unsure{remonté|remonter}}
Nb de facettes de dim $d$ $(0\leq d\unsure{\leq|<}n)$
\[
c_d=\binom nd2^{n-d}
\]
On a la formule d'EP
\[
\begin{aligned}
\sum_0^n(-1)^d c_d&=\sum(-1)^d\binom nd2^{n-d}=(-1)^n\sum\deleted{\unsure{(-1)}}\binom nd1^d(-2)^{n-d}\\
&=(-1)^n(1-2)^n=1\qquad\text{i.e.}
\end{aligned}
\]
\[
\boxed{\sum_0^n(-1)^d c_d=1}
\]
(\uline{valable en fait pour tt polyèdre convexe de dim $n$!})
% archive sheet number: 53
